\scanpage{241}
\mspage{115}
% archive sheet number: 240

$2^\circ)$ Donc le corollaire 8.5, \deleted{\unsure{la partie|la p.}} \unsure{relativement} aux conditions $b)$ $c)$ $d)$ du $1^\circ$ (à l'exclusion du $a)$, qui suppose $W$ stable par fact. directs) est valable pour un $W$ satisfaisant

\[
W(1,2\mathrm{bis},3,\underset{\inserted{\substack{\text{carré}\\\text{cocartésien}}}}{7\mathrm{bis}})
\]

$3^\circ)$ Pour \deleted{\unsure{la|la c.}} validité de 8.4. $2^\circ$, i.e. pour avoir un triple de Quillen $(W_A,C,F)$, il nous faut \inserted{(presque)} tout l'éventail des conditions sur $W$

\[
W(1\mathrm{bis},\underline{2\mathrm{bis}},3,\deleted{5},\underset{\inserted{\text{limites}}}{6},
\underset{\inserted{\substack{\text{carré}\\\text{coc.}}}}{7\mathrm{bis}},
\underset{\inserted{\text{accessibilité}}}{8})
\]

% slanted
\margin{\uline{NB} On devrait \unsure{renuméroter} les axiomes $W(\ldots.)$,}

à l'exclusion de $W(4)$, et en se \unsure{contentant}\deleted{\unsure{ant}} pour $W(2)$ de la forme affaiblie $W(2\mathrm{bis})$ sans plus. (Même l'\unsure{axiome} de connexité $W(5)$ n'a pas servi, mais \deleted{\unsure{que}} \inserted{\unsure{l'avais|au vrai|dans l'état} \unsure{d'état|actuel}} des choses, \deleted{\unsure{peut-être que}} car il y a peu de chances qu'on aura jamais à travailler avec des $W$ qui n'y satisfont pas ! \deleted{\unsure{)}} Il se pourrait d'ailleurs, en présence des autres axiomes, ou seulement de $W(1\mathrm{bis},2\mathrm{bis},3,7\mathrm{bis})$, que si $W(5)$ n'est \uline{pas} satisfait, on ait $W=\mathrm{Fl}\,\mathrm{Cat}$ tout entier, donc $W_A=\mathrm{Fl}\,\widehat A$ !)
